\section{Placement Properties and Residual-Window Analysis}
\label{sec:properties}

To move beyond a single artifact, we state the guarantees as TEE-agnostic
properties and argue the \emph{architecture} --- not the specific SEV-SNP/Azure
instantiation --- enforces them. Let a placement be the binding of a sensitive
pod $p$ to a node $n$ under policy $\pi$, using evidence $e$, producing a signed
decision $d$ with token $t$.

\begin{description}
  \item[P1 Placement Safety.] $p$ binds to $n$ only if there exists evidence
  $e$ for $n$ that is \emph{verified} (appraised by the central verifier),
  \emph{non-revoked}, and satisfies $\pi$'s TEE requirement, at the instant of
  bind.
  \item[P2 Evidence Freshness.] The $e$ used to bind satisfies
  $\mathit{age}(e) \le \pi.\mathit{maxAge}$ at the instant of bind; evidence that
  expires or is revoked between selection and bind cannot be used.
  \item[P3 Identity Binding.] $t$ is valid only for the exact tuple
  $\langle \mathrm{uid}(p), H(\mathrm{spec}(p)), n, H(e), H(\pi)\rangle$; any
  mismatch is rejected offline.
  \item[P4 Label Resistance.] No attacker-writable metadata (labels,
  \texttt{nodeSelector}) can cause a bind that P1 would forbid.
  \item[P5 Runtime Consistency.] The effective runtime class equals a
  $\pi$-allowed class, and simulated classes are refused under production mode.
  \item[P6 Verifiable Placement.] Given only the scheduler's public key and $t$,
  a third party decides whether $d$ is a genuine, unmodified decision for $p$ on
  $n$ under $\pi$ and $e$, without cluster access.
\end{description}

\paragraph{Why the architecture enforces them.}
P1/P2 hold because \emph{selection} (Filter/Score) and \emph{commitment}
(PreBind/Bind) both consult the central verifier's \texttt{AttestationEvidence}
and the two are separated by the fail-closed PreBind re-check
(\S\ref{sec:design}); no path binds without a fresh, verified, non-revoked read.
P3/P6 hold because $t$ is an EUF-CMA signature over the canonical tuple: forging
a valid $t$ for a different tuple reduces to forging Ed25519. P4 holds because the
scheduler consults evidence, never labels, for the safety decision (falsified in
E5: with the scheduler removed, a forged label bypasses; with it, it does not).
P5 is enforced at admission and re-checked at PreBind. Crucially, these arguments
mention SEV-SNP only through the abstract predicate ``$e$ is verified''; swapping
the TEE or attestation service changes only the verifier's appraisal, not the
scheduling guarantees --- the result is transferable.

\paragraph{Residual-window analysis.}
The scheduling cycle has four time-of-check points where mutable state ($\pi$,
$e$, node taints) is read: PreFilter, Filter, Score/Reserve, and PreBind. A
TOCTOU adversary (ADV-4) wins if state changes after the \emph{last} check that
gates the bind and before the bind itself. We therefore make PreBind the last
check and perform it \emph{immediately} before issuing the Binding, with no
awaited step in between: after PreBind re-reads $\pi$ and $e$ and mints $t$, the
only remaining operation is the \texttt{Binding} \texttt{Create}. The residual
window is thus the Reserve$\to$Bind local interval (creating $d$ and issuing the
Binding), during which no external actor is awaited; unlike a
\texttt{schedulingGate} design (B4), there is no gate-removal event that hands
control back to the API server before the bind. E4 measures this difference: B4's
window is the gate-removal$\to$bind interval an external revocation can hit
($\sim$1.1\,s median), whereas B5's is the in-process Reserve$\to$Bind interval
with the re-check inside it, empirically $0$ at one-second resolution. We do not
claim the residual in-process interval is provably empty --- a revocation landing
between the PreBind read and the Binding create is a theoretical race --- but it
is bounded by one API write and is not externally schedulable, which is the
qualitative gap between B5 and any check-then-ungate baseline.
